%% ma: L12-B08-0160
%% lop: 12
%% chuong: 2
%% bai: 8
%% dang: 12.8.51
%% loai: TN4
%% muc: TH
%% dap_an: "B"
%% nguon: "CĐ2.4.pdf giáo viên gửi 27/10/2026 (M07, Lê Minh Khôi – Nguyễn Anh Khôi, Chương II Chủ đề 4 Các bài toán về tâm tỉ cự), Câu 10"
%% kien_thuc: "Bài 8 (biểu thức tọa độ của phép toán vectơ, độ dài vectơ), Bài 6 (điểm thỏa hệ thức vectơ, quy tắc ba điểm)"
%% trang_thai: chinh-thuc
%% ghi_chu: "Giáo viên đã duyệt 27/10/2026"
\begin{ex}
Trong không gian $Oxyz$, cho ba điểm $A(2;-4;1)$, $B(9;3;5)$, $C(3;1;-1)$. Xét điểm $H(a;b;c)$ thay đổi trong không gian sao cho $\left|-5\overrightarrow{HA}+\overrightarrow{HB}+2\overrightarrow{HC}\right|$ đạt giá trị nhỏ nhất. Tọa độ của điểm $H$ là
\choice{$H\left(\dfrac52;\dfrac{25}2;-1\right)$}{\True $H\left(-\dfrac52;-\dfrac{25}2;1\right)$}{$H\left(-\dfrac52;\dfrac{25}2;1\right)$}{$H\left(\dfrac52;\dfrac{25}2;1\right)$}
\loigiai{Gọi $I$ thỏa $-5\overrightarrow{IA}+\overrightarrow{IB}+2\overrightarrow{IC}=\vec 0\Rightarrow I=\dfrac{-5A+B+2C}{-2}=\left(-\dfrac52;-\dfrac{25}2;1\right)$.

Biểu thức bằng $2HI$, nhỏ nhất khi $H\equiv I$.}
\end{ex}
